% % ref: Discrete Mathematics and Its Applications by Kenneth Rosen
% % ref: Combinatorial Optimization: Theory and Algorithms
% \section{Graph Theory and the Shortest Path Problem}

% \subsection{Graph Theory Fundamentals}

% \subsubsection{Basic Definitions and Concepts}

% Formally, a graph is defined as an ordered pair $G = (\mathcal{V}, \mathcal{E})$, comprising a set of vertices $\mathcal{V}$ and a set of edges $\mathcal{E}$, where each edge connects a pair of vertices. The structural classification of the graph depends on the ordering of these pairs: unordered pairs $\{u, w\}$ denote an undirected graph, whereas ordered pairs $(u, w)$ characterize a directed graph. In the context of optimization, the graph is typically augmented with a cost function $c: \mathcal{E} \rightarrow \mathbb{R}$, assigning a real-valued weight to each edge.

% Key topological properties include the vertex degree, which counts incident edges; in directed graphs, this is distinguished into in-degree ($\text{deg}^-$) and out-degree ($\text{deg}^+$). Connectivity is described via paths—finite sequences of distinct vertices and edges—where the path length is the aggregate cost of its constituent edges. A cycle is defined as a path that originates and terminates at the same vertex. Optimization problems generally operate on a search domain defined by subgraphs $H = (W, F)$, formed by subsets of vertices $W \subseteq \mathcal{V}$ and edges $F \subseteq \mathcal{E}$ that maintain incidence relationships.

% \subsubsection{Combinatorial Optimization on Graphs}

% Combinatorial optimization problems on graphs generally seek a substructure that minimizes an objective function subject to validity constraints. For a weighted graph $G = (\mathcal{V}, \mathcal{E})$ and a set of feasible subgraphs $\mathcal{H}$, the problem is to identify a subgraph $H = (W, F) \in \mathcal{H}$ that minimizes the cumulative weight of its components. This is formally expressed as:

% \begin{align}
% \text{minimize} \quad & \sum_{v \in W} c_v + \sum_{e \in F} c_e \tag{1.1a} \\
% \text{subject to} \quad & H=(W, F) \in \mathcal{H}. \tag{1.1b}
% \end{align}

% To utilize standard mathematical programming techniques, such as Branch and Bound, the problem is transformed into an Integer Linear Programming (ILP) formulation. This involves parameterizing the search space using incidence vectors $\mathbf{y}^H \in \{0, 1\}^{|\mathcal{V} \cup \mathcal{E}|}$, where binary variables indicate the inclusion of specific vertices or edges in the solution. Geometrically, the set of valid subgraphs corresponds to a polytope $\mathcal{Y}$ in Euclidean space. The intersection of this polytope with the integer grid, $\mathcal{Y} \cap \{0, 1\}^{|\mathcal{V} \cup \mathcal{E}|}$, yields the incidence vectors of feasible subgraphs. Consequently, the discrete optimization problem is recast as a linear optimization over an integer domain:

% \begin{align}
% \text{minimize} \quad & \sum_{v \in \mathcal{V}} c_v y_v + \sum_{e \in \mathcal{E}} c_e y_e \tag{1.2a} \\
% \text{subject to} \quad & \mathbf{y} \in \mathcal{Y} \cap \{0, 1\}^{|\mathcal{V} \cup \mathcal{E}|}. \tag{1.2b}
% \end{align}

% This formulation allows for the application of convex analysis, ensuring that the optimal solution to the algebraic model corresponds exactly to the optimal substructure in the original graph problem.

% \subsection{The Shortest Path Problem (SPP)}

% \subsubsection{Definition and Classical Algorithms}

% The Shortest Path Problem (SPP) on a weighted graph $G=(\mathcal{V}, \mathcal{E})$ involves identifying a path $P$ between a source $s$ and a destination $t$ such that the summation of edge weights is minimized. If $P$ comprises a sequence of edges $e_1, \ldots, e_k$, the objective is to minimize $\sum_{i=1}^{k} c_{e_i}$.

% Classical approaches to solving the SPP vary based on edge weight characteristics. Dijkstra's algorithm employs a greedy strategy suitable for graphs with non-negative weights ($c_e \ge 0$). For graphs containing negative edge weights, provided no negative cycles exist, the Bellman-Ford algorithm is utilized, albeit with higher computational complexity. The Floyd-Warshall algorithm extends this to the all-pairs shortest path problem, accommodating negative edges but remaining sensitive to negative cycles.

% \subsubsection{Network Flow Formulation for the Shortest Path Problem}

% The Shortest Path Problem on a directed graph $G = (\mathcal{V}, \mathcal{E})$ with non-negative edge costs $c_{uv} \ge 0$ can be rigorously modeled as a Minimum Cost Flow problem involving a single unit of flow. We define binary decision variables $x_{uv} \in \{0, 1\}$ for each edge $(u, v) \in \mathcal{E}$, where $x_{uv} = 1$ indicates the edge is part of the optimal path. The optimization model is formulated as follows:

% \begin{subequations}
% \begin{align}
% \text{minimize} \quad & \sum_{(u, v) \in \mathcal{E}} c_{uv} x_{uv} \label{eq:spp_obj} \\
% \text{subject to} \quad & \sum_{v \in \mathcal{V}^+(u)} x_{uv} - \sum_{v \in \mathcal{V}^-(u)} x_{vu} =
% \begin{cases}
% 1 & \text{if } u = s \\
% -1 & \text{if } u = t \\
% 0 & \text{otherwise}
% \end{cases}, \quad \forall u \in \mathcal{V} \label{eq:flow_conservation} \\
% & \sum_{v \in \mathcal{V}^-(u)} x_{vu} \le 1, \quad \forall u \in \mathcal{V} \setminus \{s, t\} \label{eq:node_capacity} \\
% & x_{uv} \ge 0, \quad \forall (u, v) \in \mathcal{E} \label{eq:non_negativity}
% \end{align}
% \end{subequations}

% Here, $\mathcal{V}^+(u)$ and $\mathcal{V}^-(u)$ denote the sets of successors and predecessors of node $u$, respectively. The constraints in this formulation enforce specific structural properties. Equation \eqref{eq:flow_conservation} represents flow conservation (Kirchhoff's law), ensuring that a net flow of one unit leaves the source $s$, enters the target $t$, and is conserved at all intermediate nodes. Equation \eqref{eq:node_capacity} imposes a node capacity constraint, restricting the inflow at any intermediate vertex to unity, thereby enforcing a simple path structure where no vertex is revisited. Given non-negative costs, cycles are inherently suboptimal, rendering explicit cycle-elimination constraints unnecessary.

% A significant advantage of this formulation lies in the algebraic properties of the constraint matrix derived from \eqref{eq:flow_conservation}. This node-arc incidence matrix exhibits Total Unimodularity (TUM). According to combinatorial optimization theory, if the constraint matrix is TUM and the right-hand side vector is integral, every basic feasible solution of the Linear Programming (LP) relaxation is guaranteed to be integral~\cite{Ahuja1993}. Consequently, despite the continuous non-negativity constraints in \eqref{eq:non_negativity}, the problem can be solved efficiently using standard LP algorithms to yield a valid binary path, obviating the need for computationally intensive integer programming methods.



\section{Lý thuyết Đồ thị và Bài toán Đường đi Ngắn nhất}

\subsection{Cơ sở Lý thuyết Đồ thị}

\subsubsection{Các định nghĩa và khái niệm cơ bản}

Một cách hình thức, một đồ thị được định nghĩa là một cặp có thứ tự
$G = (\mathcal{V}, \mathcal{E})$, bao gồm một tập các đỉnh $\mathcal{V}$ và một
tập các cạnh $\mathcal{E}$, trong đó mỗi cạnh nối một cặp đỉnh. Việc phân loại
cấu trúc của đồ thị phụ thuộc vào tính có thứ tự của các cặp này: các cặp không
có thứ tự $\{u, w\}$ biểu diễn đồ thị vô hướng, trong khi các cặp có thứ tự
$(u, w)$ đặc trưng cho đồ thị có hướng. Trong bối cảnh tối ưu hóa, đồ thị thường
được mở rộng bằng một hàm chi phí
$c: \mathcal{E} \rightarrow \mathbb{R}$, gán một trọng số thực cho mỗi cạnh.

Các tính chất tô pô quan trọng bao gồm bậc của đỉnh, là số cạnh liên thuộc với
đỉnh đó; trong đồ thị có hướng, khái niệm này được phân biệt thành bậc vào
($\text{deg}^-$) và bậc ra ($\text{deg}^+$). Tính liên thông được mô tả thông qua
các đường đi—là các dãy hữu hạn gồm các đỉnh và cạnh phân biệt—trong đó độ dài
đường đi được xác định bằng tổng chi phí của các cạnh cấu thành. Một chu trình
được định nghĩa là một đường đi bắt đầu và kết thúc tại cùng một đỉnh. Các bài
toán tối ưu thường được xét trên một miền tìm kiếm được xác định bởi các đồ thị
con $H = (W, F)$, hình thành từ các tập con của đỉnh $W \subseteq \mathcal{V}$
và các cạnh $F \subseteq \mathcal{E}$, đồng thời bảo toàn các quan hệ liên thuộc.

\subsubsection{Tối ưu hóa tổ hợp trên đồ thị}

Các bài toán tối ưu hóa tổ hợp trên đồ thị thường nhằm tìm một cấu trúc con sao
cho hàm mục tiêu đạt giá trị nhỏ nhất, đồng thời thỏa mãn các ràng buộc hợp lệ.
Với một đồ thị có trọng số $G = (\mathcal{V}, \mathcal{E})$ và một tập các đồ thị
con khả thi $\mathcal{H}$, bài toán là xác định một đồ thị con
$H = (W, F) \in \mathcal{H}$ sao cho tổng trọng số của các thành phần của nó là
nhỏ nhất. Bài toán này được biểu diễn hình thức như sau:
\begin{align}
\text{tối thiểu hóa} \quad & \sum_{v \in W} c_v + \sum_{e \in F} c_e \tag{1.1a} \\
\text{với điều kiện} \quad & H=(W, F) \in \mathcal{H}. \tag{1.1b}
\end{align}

Để có thể áp dụng các kỹ thuật lập trình toán học tiêu chuẩn, chẳng hạn như
Branch-and-Bound, bài toán được chuyển đổi sang dạng Quy hoạch Tuyến tính
Nguyên (Integer Linear Programming -- ILP). Quá trình này bao gồm việc tham số
hóa không gian tìm kiếm bằng các vectơ liên thuộc
$\mathbf{y}^H \in \{0,1\}^{|\mathcal{V} \cup \mathcal{E}|}$, trong đó các biến nhị
phân biểu thị việc một đỉnh hoặc một cạnh cụ thể có được chọn vào nghiệm hay
không. Về mặt hình học, tập các đồ thị con hợp lệ tương ứng với một đa diện
$\mathcal{Y}$ trong không gian Euclid. Giao của đa diện này với lưới số nguyên,
$\mathcal{Y} \cap \{0,1\}^{|\mathcal{V} \cup \mathcal{E}|}$, tạo thành các vectơ
liên thuộc của các đồ thị con khả thi. Do đó, bài toán tối ưu rời rạc ban đầu
được tái biểu diễn thành một bài toán tối ưu tuyến tính trên miền nguyên:
\begin{align}
\text{tối thiểu hóa} \quad &
\sum_{v \in \mathcal{V}} c_v y_v + \sum_{e \in \mathcal{E}} c_e y_e \tag{1.2a} \\
\text{với điều kiện} \quad &
\mathbf{y} \in \mathcal{Y} \cap \{0,1\}^{|\mathcal{V} \cup \mathcal{E}|}. \tag{1.2b}
\end{align}

Cách phát biểu này cho phép áp dụng các công cụ của phân tích lồi, đồng thời đảm
bảo rằng nghiệm tối ưu của mô hình đại số tương ứng chính xác với cấu trúc con
tối ưu của bài toán đồ thị ban đầu.

\subsection{Bài toán Đường đi Ngắn nhất (SPP)}

\subsubsection{Định nghĩa và các thuật toán cổ điển}

Bài toán Đường đi Ngắn nhất (Shortest Path Problem -- SPP) trên một đồ thị có
trọng số $G = (\mathcal{V}, \mathcal{E})$ nhằm xác định một đường đi $P$ từ đỉnh
nguồn $s$ đến đỉnh đích $t$ sao cho tổng trọng số các cạnh là nhỏ nhất. Nếu $P$
bao gồm một dãy các cạnh $e_1, \ldots, e_k$, mục tiêu là tối thiểu hóa
$\sum_{i=1}^{k} c_{e_i}$.

Các phương pháp cổ điển để giải SPP phụ thuộc vào đặc tính của trọng số cạnh.
Thuật toán Dijkstra sử dụng chiến lược tham lam và phù hợp với các đồ thị có
trọng số không âm ($c_e \ge 0$). Đối với các đồ thị có cạnh mang trọng số âm,
miễn là không tồn tại chu trình âm, thuật toán Bellman--Ford được sử dụng, mặc
dù có độ phức tạp tính toán cao hơn. Thuật toán Floyd--Warshall mở rộng bài toán
này cho trường hợp tìm đường đi ngắn nhất giữa mọi cặp đỉnh, cho phép tồn tại
trọng số âm nhưng vẫn nhạy cảm với sự xuất hiện của các chu trình âm.

\subsubsection{Mô hình dòng mạng cho Bài toán Đường đi Ngắn nhất}

Bài toán Đường đi Ngắn nhất trên một đồ thị có hướng
$G = (\mathcal{V}, \mathcal{E})$ với chi phí cạnh không âm $c_{uv} \ge 0$ có thể
được mô hình hóa một cách chặt chẽ như một bài toán Dòng Chi phí Tối thiểu với
một đơn vị dòng. Ta định nghĩa các biến quyết định nhị phân
$x_{uv} \in \{0,1\}$ cho mỗi cạnh $(u,v) \in \mathcal{E}$, trong đó
$x_{uv} = 1$ cho biết cạnh đó thuộc đường đi tối ưu. Mô hình tối ưu được phát
biểu như sau:
\begin{subequations}
\begin{align}
\text{minimize} \quad & \sum_{(u, v) \in \mathcal{E}} c_{uv} x_{uv} \label{eq:spp_obj} \\
\text{subject to} \quad & \sum_{v \in \mathcal{V}^+(u)} x_{uv} - \sum_{v \in \mathcal{V}^-(u)} x_{vu} =
\begin{cases}
1 & \text{if } u = s \\
-1 & \text{if } u = t \\
0 & \text{otherwise}
\end{cases}, \quad \forall u \in \mathcal{V} \label{eq:flow_conservation} \\
& \sum_{v \in \mathcal{V}^-(u)} x_{vu} \le 1, \quad \forall u \in \mathcal{V} \setminus \{s, t\} \label{eq:node_capacity} \\
& x_{uv} \ge 0, \quad \forall (u, v) \in \mathcal{E} \label{eq:non_negativity}
\end{align}
\end{subequations}

Ở đây, $\mathcal{V}^+(u)$ và $\mathcal{V}^-(u)$ lần lượt ký hiệu tập các đỉnh kế
tiếp (successors) và các đỉnh tiền nhiệm (predecessors) của đỉnh $u$. Các ràng
buộc trong mô hình này đảm bảo những tính chất cấu trúc cụ thể. Phương trình
\eqref{eq:flow_conservation} biểu diễn điều kiện bảo toàn dòng (định luật
Kirchhoff), đảm bảo rằng một đơn vị dòng rời khỏi đỉnh nguồn $s$, đi vào đỉnh
đích $t$, và được bảo toàn tại tất cả các đỉnh trung gian. Ràng buộc
\eqref{eq:node_capacity} áp đặt giới hạn dung lượng nút, giới hạn dòng đi vào
mỗi đỉnh trung gian không vượt quá một đơn vị, qua đó đảm bảo cấu trúc đường đi
đơn, trong đó không có đỉnh nào bị ghé thăm nhiều hơn một lần. Với chi phí
không âm, các chu trình luôn là không tối ưu, do đó không cần thiết phải bổ sung
các ràng buộc loại bỏ chu trình một cách tường minh.

Một ưu điểm quan trọng của cách mô hình hóa này nằm ở các tính chất đại số của
ma trận ràng buộc thu được từ~\eqref{eq:flow_conservation}. Ma trận liên thuộc
đỉnh-cạnh này có tính \textit{hoàn toàn đơn mô} (Total Unimodularity--TUM).
Theo lý thuyết tối ưu hóa tổ hợp, nếu ma trận ràng buộc là TUM và vectơ vế phải
là nguyên, thì mọi nghiệm khả thi cơ sở của bài toán thư giãn Quy hoạch Tuyến
tính (LP) đều được đảm bảo là nguyên~\cite{AhujaMagnantiOrlin1993}. Do đó, mặc dù chỉ áp đặt
các ràng buộc không âm liên tục trong~\eqref{eq:non_negativity}, bài toán vẫn có
thể được giải hiệu quả bằng các thuật toán LP tiêu chuẩn để thu được một đường
đi nhị phân hợp lệ, mà không cần đến các phương pháp lập trình nguyên có chi phí
tính toán cao.
